\section{How do experience and signal reliability change the perceived meaning of a sound?}
\label{sec:q5}

% Part III — Updating responses over time (Q5).
% See docs/OUTLINE.md — Q5.
% Draft 1. Style model: Magrath (agents/STYLE.md).

Associative learning can give a sound defensive meaning, as experience does for rat 22-kHz calls (\secref{sec:q2}). A receiver can also learn something less obvious: how far a signal can be trusted.
% MISSING: callers differ in age, experience, and tendency to call at harmless events; suggested source: BlumsteinDaniel2004 covers age class only; search sciurid false-alarm rates by caller age.
Alarm calls are produced by individuals that differ in age, experience, and their tendency to call at harmless events, and a receiver that tracks these differences can respond strongly to some callers and weakly to others. Reliability learning and generic habituation both reduce responses to repeated calls, but they make different predictions. Reliability learning predicts a decline specific to a caller or context, and selective renewal when the caller changes or a call is again followed by danger; habituation predicts a decline that generalizes across callers of similar acoustic structure and recovers with time. This section asks which of these processes the evidence supports and in which rodents.

Discriminating callers is a precondition for learning their reliability, and it is now documented beyond sciurids. Juvenile Richardson's ground squirrels hearing repeated calls from one individual reduced their vigilance by the second call; a new exemplar from the same caller did not restore vigilance, whereas a call from a different individual did \citep{Hare1998}. Yellow-bellied marmots likewise habituated to repeated calls from one individual and recovered when caller identity changed \citep{BlumsteinDaniel2004}. In laboratory rats, individuals preferred to self-administer playback of 50-kHz calls from strangers over those of cagemates, only stranger calls reduced cocaine intake, and lesions of the subthalamic nucleus abolished both effects \citep{VielleEtAl2021}; some stranger exemplars were repeated, which the authors acknowledge as a limitation.
% VERIFY: WarrenEtAl2025 (not full-text checked): partner versus stranger discrimination, including calls directed at another female.
Female prairie voles discriminated the calls of their partner from those of a stranger male, including calls the partner produced toward another female \citep{WarrenEtAl2025}.
% VERIFY: HareWarkentin2012 (not full-text checked): no differential response to mother's or colony members' calls.
Discrimination does not mean that every social category changes behavior: juvenile Richardson's ground squirrels did not respond differently to their mother's calls or to calls from members of their own colony \citep{HareWarkentin2012}, and marmots did not distinguish familiar from unfamiliar callers \citep{BlumsteinDaniel2004}.
Identity is therefore widely available to rodent receivers, but familiarity and kinship are not, by themselves, reasons to respond more strongly.

Direct tests of reliability learning come only from sciurids, and they point in opposite directions. \citet{HareAtkins2001} presented juvenile female Richardson's ground squirrels with ten exemplars each from two callers: calls from the reliable caller were usually paired with a stuffed badger, whereas calls from the unreliable caller were never accompanied by a predator. Postural responses to the unreliable caller declined and differed from responses to the reliable caller by the fifth and tenth playbacks; vigilance duration followed the same pattern only as a nonsignificant trend, and samples were small. In yellow-bellied marmots, responses depended on caller reliability, on simulated distance, and on whether one or two individuals called, but marmots responded more, not less, to calls from unreliable individuals \citep{BlumsteinEtAl2004}. Blumstein and colleagues interpret the increase as investment in independent assessment of risk when a caller's alarms cannot be taken at face value. The two studies differ in species, receiver age, how reliability was manipulated, and endpoint (posture and vigilance in ground squirrels, foraging suppression and vigilance in marmots), and any of these could contribute to the difference in direction.

Error management offers a way to treat the divergence as a prediction rather than a contradiction.
% VERIFY: Nesse2005 (not full-text checked): whether the smoke-detector argument covers callers with a history of false alarms and cheaper information sources, or whether this is our extension.
When missing a real threat is costly, a receiver should keep responding to a caller with a history of false alarms unless it has a cheaper source of information \citep{Nesse2005}. Discounting an unreliable caller should be favored when independent checking is costly or ineffective, and scrutiny when the receiver can assess the threat itself at low cost, for example by scanning from a safe position. This account predicts that the direction of reliability effects should vary with the cost of independent assessment, which could differ with age, habitat visibility, and distance from refuge. It has not been tested, and the existing data do not distinguish it from simpler explanations, such as differences in how strongly each design established a caller's unreliability.
% VERIFY: LillyEtAl2019 (not full-text checked): vigilance relaxation with heterospecific chatter.
Updating can also draw on evidence that danger has passed, independently of the alarm signal's history, as when gray squirrels use heterospecific chatter as a safety cue \citep{LillyEtAl2019} (\secref{sec:q2}).

Laboratory rodent studies have shown the components of this process separately but not together. Rats and prairie voles discriminate individuals by their calls \citep{VielleEtAl2021,WarrenEtAl2025}; experience gives 22-kHz calls defensive meaning \citep{KimEtAl2010,CalubEtAl2018}; and maternal experience reshapes the sensory representation of pup calls \citep{SchiavoEtAl2020}. We found no laboratory study in which the reliability of a caller was manipulated, as in the sciurid experiments, and responses to that caller were tracked. The bridge between field and laboratory systems is therefore reliability, not identity. Whether a declining response reflects learned unreliability or generic habituation can be decided only by designs that vary the caller while holding acoustic structure and schedule constant, and by testing whether renewal is specific to the caller whose reliability changed.

\tabref{tab:summary-q5} consolidates these claims for later condensation or figure drafting.

% Intermediate summary matrix for Q5. May later be condensed, moved to an appendix, or redrawn as a figure.
\begin{table}[tbp]
\centering
\footnotesize
\singlespacing
\caption{Working summary for Q5: identity versus reliability. Distinguishing prediction: reliability learning is caller- or context-specific and renews selectively; habituation generalizes across similar callers and recovers with time.}
\label{tab:summary-q5}
\begin{tabular}{@{}p{1.3in}p{2.4in}p{1.4in}p{0.9in}@{}}
\toprule
\textbf{Process} & \textbf{Finding} & \textbf{Where shown} & \textbf{Status} \\
\midrule
Caller identity & Discriminated; renewal when caller changes & Sciurids; rats; prairie voles & Strong \\
\addlinespace[2pt]
Familiarity / kinship & Often no behavioral effect alone & Marmots (familiarity null); ground squirrels (mother/colony null) & Nulls matter \\
\addlinespace[2pt]
Reliability: discount & Responses fall to callers unpaired with predators & Juvenile Richardson's ground squirrels & Field only \\
\addlinespace[2pt]
Reliability: scrutinize & Responses rise to unreliable callers & Yellow-bellied marmots & Field only \\
\addlinespace[2pt]
Direction of reliability & May depend on cost of independent assessment & Prediction; untested & Open \\
\addlinespace[2pt]
Safety updating & Heterospecific chatter speeds relaxation after hawk call & Gray squirrels & Supporting \\
\addlinespace[2pt]
Lab reliability test & Identity present; reliability manipulated as in sciurids & --- & Gap \\
\bottomrule
\end{tabular}
\end{table}

